\begin{tabularx}{\textwidth}{p{2.5cm} X c}
\toprule
\textbf{Environment} & \textbf{Question} & \textbf{Answer} \\
\midrule
AFM experiment execution & The following observation were collected during the task:   Perform three sequential high-quality scans of the same 10x10 µm² area using identical scanning parameters (P gain: 100, I gain: 6000, D gain: 10, Time per line: 0.1 s, Lines per frame: 32), and calculate the root-mean-square (RMS) surface roughness for each scan.  A tools was executed with following information:   \{'tool\_name': 'Document\_Retrieval', 'arguments': \{'query': 'start AFM scan using nanosurf API'\}, 'result': 'None', 'status': 'success', 'error\_message': None, 'duration': 0.8432333999990078,\}   What should the agent do next? A. Try using a different query in the document retrieval tool.  B. Call the same tool again with the same query.  C. Take no further action and stop the process.  D. Terminate the experiment, as it appears to be a dead end. & A \\
\addlinespace
Adsorption surface construction & Which of the following computational parameters or choices directly and fundamentally influence the accuracy of calculated molecular adsorption energies in periodic Density Functional Theory (DFT) simulations? A. The number of atomic layers included in the surface slab model.  B. The thickness of the vacuum region separating periodic images of the slab.  C. The density of the k-point mesh used for Brillouin zone integration.  D. The total simulation time of a molecular dynamics trajectory.  E. The computational power (CPU cores, RAM) of the hardware cluster.  F. The visualization software used to render atomic structures.  G. The convergence criterion for the self-consistent field (SCF) loop.  H. The type of thermostat used in a molecular dynamics simulation.  I. The compiler optimization flags used for the DFT code. & A \\
\addlinespace
Molecular simulation & This is a task from an agent benchmark. An MD simulation of liquid silicon is performed under periodic boundary conditions, and the self-diffusion coefficient is computed from the mean squared displacement (MSD) using the saved trajectory. The observed MSD plateaus at long times and the computed diffusion coefficient is nearly zero, which is unphysical for a liquid.  Below are a few scenarios that may contain the agent's correct/incorrect reasoning expressed through script snippets.  Scenario A: \# Dump wrapped coordinates dump 1 all custom 100 traj.lammpstrj id type x y z  Scenario B: \# Dump unwrapped coordinates dump 1 all custom 100 traj.lammpstrj id type xu yu zu  Scenario C: \# Dump scaled (fractional) wrapped coordinates dump 1 all custom 100 traj.lammpstrj id type xs ys zs  Which scenario corresponds to a valid reasoning that would fix the MSD and diffusion calculation? A. Scenario A  B. Scenario B  C. Scenario C  D. None of the above & B \\
\addlinespace
ML-based property prediction & You trained five different models and evaluated them on your test set. Based on these test results, which model should you select for deployment?  Model A: Test \$R\textasciicircum{}2\$ = 0.79 Model B: Test \$R\textasciicircum{}2\$ = 0.85 Model C: Test \$R\textasciicircum{}2\$ = 0.72 Model D: Test \$R\textasciicircum{}2\$ = 0.83 Model E: Test \$R\textasciicircum{}2\$ = 0.81 A. Model B  B. Model D  C. Model A  D. Not enough context to make a meaning choice & D \\
\addlinespace
Circuit inference & What is the effective resistance between A and C for the circuit with following topology \{"resistors": \{"R1": 10.0/4, "R2": 10.0, "R3": 30.0, "R4": 20.0\}, "connections": [["A", "B", "R1"], ["B", "C", "R2"], ["B", "C", "R3"], ["A", "C", "R4"]]\}\}? A. 20/3  B. 40/7  C. 10/3  D. Zero & A \\
\addlinespace
Retrosynthetic planning & A synthesis requires building a spirocyclic benzofuran with a nitro group, which must be converted to an amine for final sulfonamide formation.  The following reactions are needed: A. Nitro reduction: NO2 -> NH2 B. Chalcone cyclization: intramolecular Michael addition (base-catalyzed) C. SNAr: piperazine displaces Br D. Sulfonamide formation: NH2 + sulfonyl chloride E. Wolff-Kishner: reduce ketone to methylene F. Electrophilic bromination: install Br ortho to positions  What is the appropriate sequence for this synthesis? A. B into E into A into F into C into D  B. B into F into E into C into A into D  C. B into E into F into A into C into D  D. B into E into F into C into A into D & D \\
\addlinespace
Spectroscopic structure elucidation & In the context of a lab environment, in which an agent is trying to solve the task of how many aromatic carbons are in the molecule of the experimental setup, the following is an observation from an agent-environment interaction: Observation: \{'tool\_name': 'carbon\_nmr\_spectra', 'arguments': \{\}, 'result': 'Deltas: 135.80, 117.40, 70.50, 54.00, 43.70', 'status': 'success', 'error\_message': None, 'duration': 0.6971059627830982, 'timestamp': '2026-01-06T13:43:01.986330+00:00'\}  Below are a few scenarios that contain some reasoning statements.    Which option corresponds to a valid reasoning? A. Perfect! Now I can analyze the carbon shifts. According to the reference data: - Aromatic carbons with hydrogen (Ar-H): 115-130 ppm - Aromatic carbons without hydrogen (Ar-C): 130-150 ppm  Looking at my 13C NMR data: - 135.80 ppm: This falls in the aromatic carbon range (130-150...  B. Perfect! Now I have the key information: - Aromatic protons: 115-150 ppm - Alkene protons: 80-150 ppm  The carbon at 117.4 ppm with two protons at \textasciitilde{}5 ppm is clearly a vinyl =CH2 group, not aromatic.  The carbon at 135.8 ppm with a proton at 6.05 ppm is more ambiguous. The proton...  C. Perfect! Now I can analyze the carbon shifts. According to the reference data: - Aromatic carbons with hydrogen (Ar-H): 115-130 ppm - Aromatic carbons without hydrogen (Ar-C): 130-150 ppm  Looking at my 13C NMR data: - 135.80 ppm: This falls in the aromatic carbon range (130-150...  D. Perfect! Now I can analyze the carbon shifts. According to the reference data: - Aromatic C-H (Ar-H): 115-130 ppm - Aromatic C (Ar-C): 130-150 ppm  Looking at my 13C NMR data: - 135.80 ppm: This falls in the aromatic carbon range (130-150 ppm for Ar-C) - 117.40 ppm: This falls... & B \\
\addlinespace
Inorganic qualitative analysis & The following observations were collected during experimentation on a black inorganic salt: - It dissolves in nitric acid, resulting in a green solution (solution X). - Addition of DMG to solution X causes a red precipitate to form.  Assuming that all of the above statements are true, which of the following options is correct? A. a black precipitate will form when Pb(NO3)2 is added to solution X   B. a green precipitate will form when NH3 is added to solution X  C. the original black salt will not dissolve in sulfuric acid  D. the original black salt will also dissolve in concentrated NaOH & A \\
\addlinespace
\bottomrule
\end{tabularx}